% !TEX root = ../../../main.tex
% 所属：第一部分：单变量分析学 / 第1章：实数理论——从Peano公理到实数的连续性
% 节：1.4 等价关系，商集与$\Z$整数的构建
% 正文按原文件顺序保存；由 ../chapter1.tex 调用。
\section{等价关系，商集与$\Z$整数的构建}
\begin{definition}
	定义一个关系是等价关系$\sim$若它满足如下条件
	\begin{enumerate}
		\item 自反性：$A \sim A$
		\item 对称性：若 $A \sim B$ 则 $B \sim A$
		\item 传递性：若 $A \sim B$ 且 $B \sim C$ 则 $A \sim C$
	\end{enumerate}
	
	
\end{definition}
\begin{definition}
	在一个集合$X$上定义一个等价关系$\sim$.
	对某一元素 $a, \left[a\right]$ 定义成所有与 $a$ 等价的元素构成的集合，也就是说
	
	$$
	[a]=\left\{x \in X: x \sim a\right\}
	$$
	
	$\left[a\right]$被称作 $a$ 的等价类.
\end{definition}

\begin{lemma}
	等价类 $[a], a \in X$ ，$[a]$ 有以下三条性质.
	\begin{enumerate}
		\item $a \in[a]$
		\item $b \in[a]$ 则 $[a]=[b]$
		\item 任意两个等价类 $[a],[b]$ 要么 $[a]=[b]$ 要么 $[a]\cap[b]=\varnothing$
	\end{enumerate}
\end{lemma}
\begin{proof}
	\begin{enumerate}
		\item $a \sim a$
		\item 任取 $x\in[a]$，即 $x\sim a$. 由 $b\in[a]$ 知 $b\sim a$，由对称性 $a\sim b$，传递性得 $x\sim b$，故 $x\in[b]$，即 $[a]\subseteq[b]$. 同理 $[b]\subseteq[a]$，故 $[a]=[b]$
		\item 若 $[a] \cap[b] \neq \varnothing$ ,设 $x \in[a] \cap[b]$ 则 $a\sim x,b\sim x,a\sim b$, 
		$[a]=[b]$
	\end{enumerate}
\end{proof}
\begin{definition}
	＂$\sim$＂是 $X$ 上的一个等价关系.由＂$\sim$＂确定的等价类构成的集合，称为 $X$ 关于 $\sim$ 的商集，记为 $X / \sim$
	
	$$
	X / \sim=\{[a]: a \in X\}
	$$
\end{definition}
现在，利用商集的概念，我们可以从自然数上构造整数.
\begin{construction}整数的构造
	$$
	\text { 令 } S=\{(a,b): a,b \in \mathbb{N}\}
	$$
	
	（虽然我们尚未定义＂减法＂，为了方便理解 $(a, b)\Leftrightarrow a-b$）\\
	在 $S$ 上定义$\sim$
	\[(a, b) \sim(c, d) \quad\text{iff}\quad a+d=b+c\]	
	
	$$
	[a+d=b+c \Leftrightarrow a-b=c-d]
	$$
	这样
	
	$$
	\Z:=S/\sim
	$$
	每个等价类 $[a, b]$ 是一个整数.
\end{construction}   



\begin{proof}
	现验证$\sim$是等价关系.
	\begin{enumerate}
		\item  自反性 $a+b=b+a \quad(a, b) \sim(a, b)$
		\item 对称性 $(a, b) \sim(c, d) \Leftrightarrow a+d=b+c \Leftrightarrow c+b=d+a \Leftrightarrow (c, d) \sim(a, b)$
		\item 
		传递性 $(a, b) \sim(c, d), ( c, d) \sim(e, f) \Leftrightarrow$
		$$
		a+d=b+c, c+f=d+e \Leftrightarrow a+f=b+e \Leftrightarrow(a, b) \sim(e, f)
		$$
	\end{enumerate}
\end{proof}

现对 $\Z$ 上的运算进行定义.\\
\begin{definition}
	加法 $[(a, b)]+[(c, d)]:=[(a+c, b+d)]$\\
	加法逆元 $-[(a, b)]:=[(b, a)]$\\
	减法 $[(a, b)]-[(c, d)]:=[(a+d, b+c)]$\\
	乘法 $[(a, b)] \cdot[(c, d)]:=[(a c+b d, a d+b c)]$\\
	通常情况下 $[(n, 0)]$ 记作 $n \quad \Rightarrow \mathbb{N} \subset \mathbb{Z}$\\
	$[(0, n)]$ 记作 $-n$\\
	$[(a,b)]$记作$a-b$\\
\end{definition}

这样，我们就严格的构造出了整数集，并且定义了其上的运算.\\
$(\Z,+)$ 是阿贝尔群，
$(\mathbb{Z}, +, \cdot)$ 是交换环\\
\begin{example}
	-2 对应 $[(3,5)],[(1,3)]$ 等,
	-4 对应 $[(1,5)],[(6,10)]$ 等
\end{example}


